\scanpage{91}
\mspage{90}
\[
(8)\qquad \operatorname{Hom}_{C'}(x,y)=\operatorname{Hom}_{\Pi_1(P)}(i_E(x),i_E(y))
\qquad x,y\in E
\]
On voit par transport de structure que
$\mu$ opère sur $C'$, de façon que
\[
C'\longrightarrow\Pi_1(P)
\]
soit compatible avec les op. de $\mu$.
Considérons l'application
\[
(9)\qquad E_0=E/\mu\xrightarrow{i_{E_0}}T\simeq P/U
\]
induite par passage aux quotients, et soit
\deleted{\unsure{donc}} $C$ le groupoïde sur $E_0$ comme
ens. d'objets, induit par $i_{E_0}$ :
\[
(10)\qquad \operatorname{Hom}_C(x_0,y_0)=
\operatorname{Hom}_{\Pi_1(T)}(i_{E_0}(x_0),i_{E_0}(y_0)).
\]
On trouve alors un diagramme
commutatif d'hom de groupoïdes
\diagram{
\begin{tikzcd}[ampersand replacement=\&]
C' \arrow[r,"{\deleted{\unsure{équivalence}}}"] \arrow[d,"f"',"{\deleted{\unsure{app. aux fibres}}}"] \& \Pi_1(P) \arrow[d] \\
C \arrow[r] \& \Pi_1(T)
\end{tikzcd}
}\quad (11)
\margin{les flèches horizontales sont des équivalences,
si $\unsure{E}\longrightarrow\pi_0(T)$ surjectif\ldots}
compatible avec les op. de $\mu$, en faisant
